\section{Related Work}
\label{sec:related}

Automated SOC alert triage sits at the intersection of machine learning-based
anomaly detection, natural language-driven autonomous agents, and operational
cybersecurity practice. We survey these three dimensions and identify the research
gaps motivating our design.

\subsection{Machine Learning for Alert Prioritization}

Early large-scale automation of SOC workflows deployed tabular statistical
classifiers over structured endpoint and network log features.
Freitas et al.~\cite{freitas2024cgr} introduced Copilot Guided Response (CGR),
deploying Gradient Boosted Decision Trees (GBDT) with Shapley Additive Explanations
(SHAP) across Microsoft enterprise tenants, achieving a macro-F1 score above 0.87
on the GUIDE dataset. The key advantage of GBDT-based approaches lies in their
calibrated probabilistic outputs and computational efficiency ($<$0.04\,s per
incident), making them suitable for high-throughput production environments.
However, CGR and similar statistical models lack semantic comprehension of
kill-chain attack progressions: they classify individual alert feature vectors
without reasoning across correlated multi-step activity sequences.

Garcia et al.~\cite{garcia2025rosca} proposed ROSCA, integrating alert correlation
with MITRE ATT\&CK technique mapping to rank alert severity. ROSCA demonstrated
scalability advantages over graph-based correlation engines but required manually
curated MITRE mappings for each deployment environment.
Freitas et al.~\cite{freitas2026aip} subsequently introduced Adaptive Incident
Prioritization (AIP), dynamically re-ranking queued incidents based on evolving
threat context, reducing analyst queue overflow by approximately 23\% in production.
Ndichu et al.~\cite{ndichu2026survey} provide a comprehensive survey of AI-driven
alert screening approaches, cataloguing 47 methods and identifying that high
false-positive tolerance remains the predominant unresolved challenge across systems.

\subsection{LLM Agents for Autonomous Cyber Investigation}

The emergence of tool-augmented LLMs enabled a qualitative shift from passive
classification to active evidence gathering. Wei et al.~\cite{wei2025cortex}
proposed CORTEX, a multi-agent debate architecture in which independent LLM
instances challenge each other's threat hypotheses before converging on a verdict.
CORTEX improved investigation consensus quality but introduced substantial per-alert
API costs (approximately \$24 per complex multi-hop investigation), limiting
deployment to Tier-2 escalation workflows.

Tsegai et al.~\cite{tsegai2026aida} provide the most rigorous empirical characterization
of this vulnerability through AIDA, evaluating five distinct tool-using agent architectures.
All five strategies exhibited a true-positive false negative rate of at least 40.4\%, with
the dominant failure mode being premature termination when initial log queries return
empty result sets~\cite{tsegai2026aida_missing}. This failure mode is structurally absent
from pure-classifier approaches and represents the critical gap exploited in our Evidence Gate design.

\subsection{Benchmark Evaluation Frameworks}

Realistic evaluation of autonomous cyber investigation agents requires multi-hop log
environments reflecting authentic enterprise attack chains. Wu et al.~\cite{wu2026excytin}
introduce ExCyTIn-Bench, comprising 57 relational log tables and 8 complete attack
chains sourced from real-world breach telemetry. The state-of-the-art investigation
reward plateaus at 0.606, indicating that existing agents correctly trace only
approximately 60\% of attack activities even under ideal laboratory conditions.

\subsection{Research Gap and Positioning}

Table~\ref{tab:literature} summarizes our systematic comparison. The analysis reveals
a consistent tripartite gap: (1) statistical classifiers lack kill-chain semantic reasoning;
(2) LLM agents suffer from the empty-query premature termination failure mode; and
(3) no existing system jointly addresses probabilistic calibration, structured knowledge
retrieval, bounded investigation planning, and empirically validated escalation thresholds.

\begin{table}[t]
\caption{Systematic Comparison of Related SOC Triage Literature}
\label{tab:literature}
\centering
\begin{tabular}{p{1.8cm}p{0.6cm}p{2.5cm}p{3.2cm}}
\toprule
\textbf{System} & \textbf{Yr} & \textbf{Core Approach} & \textbf{Key Limitation} \\
\midrule
CGR~\cite{freitas2024cgr}        & 2024 & GBDT + SHAP        & No kill-chain semantics \\
CORTEX~\cite{wei2025cortex}      & 2025 & Multi-LLM Debate   & \$24/alert API cost \\
AIDA~\cite{tsegai2026aida}       & 2026 & Tool ReAct         & $\geq$40.4\% FNR \\
ExCyTIn~\cite{wu2026excytin}     & 2026 & Multi-hop Logs     & Reward plateau 0.606 \\
Vigilance~\cite{tilbury2026vigilance} & 2026 & SOC Field Study & Complacency trap \\
Ours                             & 2026 & GBDT+RAG+ReAct+Evidence Gate & This work \\
\bottomrule
\end{tabular}
\end{table}
